\section{Related work}
\label{sec:related}
Matching a lesion across visits is a prerequisite for measuring growth and response in follow-up imaging.
Location Registration and Recognition (LRR)~\cite{sofka2010} is the closest design to ours: it generates local registration hypotheses around a supplied nodule position, refines them, and verifies them with a lesion model.
Rafael-Palou \etal~\cite{rafaelpalou2021} instead re-identify detected nodules with a 3D Siamese network, avoiding explicit registration, while Deep Lesion Tracker~\cite{cai2021} and Transformer Lesion Tracker~\cite{tang2022tlt} learn to combine lesion appearance with anatomical context and predict the earlier location directly.
These systems differ mainly in their scoring function; how candidate locations are allocated before scoring has received less attention, and it is the question we study.
Registration offers a second route to the earlier location.
Free-form deformation with B-splines~\cite{rueckert1999}, often called B-spline registration, is a standard way to estimate a continuous map between two CT scans, and discrete registration evaluates a set of candidate displacements per control point and optimizes them jointly~\cite{glocker2008dense,heinrich2013mrf}, which is close in spirit to scoring a lattice of translations.
Learned methods predict deformation fields in a single forward pass~\cite{balakrishnan2019voxelmorph,ghazipour2021}, and recent foundation models and benchmarks cover many anatomies, including lung CT~\cite{tian2024unigradicon,hering2023learn2reg}.
Studies of registration uncertainty show that good image agreement does not guarantee a reliable mapped location~\cite{eppenhof2018,luo2017}, which motivates keeping several hypotheses rather than committing to one; we use a lesion-inclusive B-spline as the registration comparator.
A third line of work learns a descriptor of anatomical position for every voxel.
Self-supervised anatomical embeddings (SAM)~\cite{yan2022sam}\footnote{Not to be confused with the Segment Anything Model.} are trained so that the same anatomical point receives similar descriptors in different scans, and a point marked in one scan is matched by finding the most similar descriptor in the other; diffusion-guided position encodings~\cite{zhuang2026diffusion} follow the same principle, localizing globally by embedding similarity and then refining locally without a prior registration.
These methods replace an explicit search over shifts with a learned similarity, but they still report a single most-similar location, so when several locations look alike the error lies in selection rather than in coverage (\cref{sec:transfer}).
We use SAM as a baseline because its released model accepts a point query directly (\cref{sec:evaluation}). Selecting a few good hypotheses from many is a recurring problem across vision.
Non-maximum suppression greedily removes detections that overlap a higher-scoring one~\cite{neubeck2006nms}, and diverse $M$-best methods return several high-scoring but mutually different solutions so that a later stage can choose among them~\cite{batra2012diverse}.
Farthest-point sampling spreads samples to cover a domain with a bounded gap~\cite{eldar1997farthest}, and randomized correspondence search propagates good matches while sampling widely~\cite{barnes2009patchmatch}.
We apply minimum-distance selection to lesion correspondence hypotheses ordered by anatomical fit, and we measure its effect on coverage separately from its effect on the final choice.
